\uuid{SPVE}
\titre{Système linéaire à paramètres/Méthode du Pivot de Gauss}
\auteur{hanani ; mzari}
\organisation{exo7}
\licence{CC BY-NC-SA 4.0}
\datecreate{2020-01-20}
\niveau{L1}
\module{Algèbre}
\chapitre{Déterminant, système linéaire}
\sousChapitre{Système linéaire, rang}
\difficulte{3}
\theme{Systèmes d'équations linéaires, Système linéaire à paramètres, Méthode du Pivot de Gauss}
\exo7qcmid{lille-2.systemes.017}

\contenu{
\begin{question}
\qclassification{200.03}
\qtags{motcle=Système linéaire à paramètres/Méthode du Pivot de Gauss}
On considère le système d'équations, d'inconnue $(x,y,z,t)\in \Rr^4$  et de paramètres des réels $m$ et $a$ :
$$(\mathtt{S})
\left\{\begin{array}{rcc}
x+y+z+mt&=&1\\
x+y+mz+t&=&a\\
x+my+z+t&=&a^2\\
mx+y+z+t&=&a^3.\end{array}\right.$$
Quelles sont les assertions vraies ?
\begin{answers}
\good{$(\mathtt{S}) \Leftrightarrow
\left\{\begin{array}{rcc}
x+y+z+mt&=&1\\
(m-1)y+(1-m)t&=&a^2-1\\
(m-1)z+(1-m)t&=&a-1\\
(1-m)(3+m)t&=&a^3+a^2+a-m-2. \\
\end{array}\right. $}
\bad{Si $m =1$,  $(\mathtt{S})$ admet une infinité de solutions.}
\bad{Si $m\neq 1$, $(\mathtt{S})$ admet une unique solution.}
\good{Si $m=-3$ et $a\neq -1$, $(\mathtt{S})$ n'admet pas de solution.}
\end{answers}
\begin{explanations}
$$(\mathtt{S}) \Leftrightarrow
\left\{\begin{array}{rcc}
x+y+z+mt&=&1\\
(m-1)y+(1-m)t&=&a^2-1\\
(m-1)z+(1-m)t&=&a-1 \\
(1-m)y+(1-m)z+(1-m^2)t&=&a^3-m\\
\end{array}\right.$$
$$ \Leftrightarrow
\left\{\begin{array}{rcc}
x+y+z+mt&=&1\\
(m-1)y+(1-m)t&=&a^2-1\\
(m-1)z+(1-m)t&=&a-1 \\
(1-m)(3+m)t&=&a^3+a^2+a-m-2.\\
\end{array}\right. $$
\begin{enumerate}
\item[-]Si $m=1$ et $a=1$,  $(\mathtt{S})  \Leftrightarrow  x+y+z+t=1 $ admet une infinité de solutions.
\item[-]Si $m=1$ et $a\neq 1$,  $(\mathtt{S})$ n'admet pas de solution.
\item[-]Si $m=-3$ et $a=-1$,  $(\mathtt{S})$ admet une infinité de solutions.
\item[-]Si $m=-3$ et $a\neq -1$,  $(\mathtt{S})$ n'admet pas de solution.
\item[-]Si $m\neq 1$ et $m\neq -3$, $(\mathtt{S})$ admet une unique solution.
\end{enumerate}
\end{explanations}
\end{question}
}
